\begin{corollary}[Endpoint weak-type $(1,1)$ bound for the complex Gaussian Riesz transform]
Let $\mathcal{M}(\gamma_n;\mathbb{C}^n)$ denote the space of
$\gamma_n$-almost-everywhere equivalence classes of measurable
$\mathbb{C}^n$-valued functions on $\mathbb{R}^n$.  There is exactly one
$\mathbb{C}$-linear map
\[
  T_{\mathbb{C}} : L^1(\gamma_n;\mathbb{C}) \longrightarrow
    \mathcal{M}(\gamma_n;\mathbb{C}^n)
\]
with the following two properties:
\begin{enumerate}
\item $T_{\mathbb{C}}$ extends the complex transform of
\noderef{ComplexRieszTransform} from $L^2$, i.e.
$T_{\mathbb{C}}(\iota_{\mathbb{C}} v) = R_\gamma^{\mathbb{C}} v$ for every
$v \in L^2(\gamma_n;\mathbb{C})$, with $\iota_{\mathbb{C}}$ the inclusion of
\noderef{ComplexToL1}; and
\item $T_{\mathbb{C}}$ satisfies the weak-type bound
\[
  \sup_{t>0}\; t\,\gamma_n\bigl(\{x :
    |T_{\mathbb{C}}f(x)|_{\mathbb{C}^n} > t\}\bigr)
    \;\leq\; \frac{8}{\pi}\,\|f\|_{L^1(\gamma_n;\mathbb{C})}
  \qquad\text{for every } f \in L^1(\gamma_n;\mathbb{C}) .
\]
\end{enumerate}
The constant $8/\pi$ does not depend on $n$.

This is the reference paper's complex-valued corollary: its displayed estimate
is clause (2), and the phrase ``the complex-linear extension'' naming the
operator it is about is clauses (1) and the $\mathbb{C}$-linearity together,
which the last sentence of its proof supplies as a \emph{unique continuous}
complex-linear extension to $L^1$.  Property (2) is precisely boundedness from
$L^1(\gamma_n;\mathbb{C})$ into $L^{1,\infty}(\gamma_n;\mathbb{C}^n)$ with
quasi-norm at most $8/\pi$, hence continuity of $T_{\mathbb{C}}$ for the weak
$L^1$ quasi-norm.  Both clauses are needed for the uniqueness assertion to have
content: clause (2) alone is satisfied by the zero map, so it is the agreement
with $R_\gamma^{\mathbb{C}}$ on $L^2$ in clause (1) that pins the operator down,
and the bound in (2) that propagates that agreement from the $L^2$ functions to
all of $L^1$.  Uniqueness is asserted in the larger class of all
$\mathbb{C}$-linear maps into $\mathcal{M}(\gamma_n;\mathbb{C}^n)$ satisfying (1)
and (2), not merely among continuous ones.
\end{corollary}
\begin{proof}
\emph{Existence.}
\noderef{ComplexL1Extension} supplies a $\mathbb{C}$-linear map
$T_{\mathbb{C}} : L^1(\gamma_n;\mathbb{C}) \to \mathcal{M}(\gamma_n;\mathbb{C}^n)$
satisfying (1) and (2); the supremum form of (2) in the present statement is the
same assertion as the per-level form there, since a family of real numbers is
bounded by a constant exactly when its supremum is.

\emph{Uniqueness.}
Let $T$ and $T'$ both be $\mathbb{C}$-linear and satisfy (1) and (2), and fix
$f \in L^1(\gamma_n;\mathbb{C})$.

\emph{Step 1: an approximating sequence.}
By clause (2) of \noderef{ComplexRieszL1Limit} applied to $f$ there is a
sequence $(v_k)$ in $L^2(\gamma_n;\mathbb{C})$ with
$\iota_{\mathbb{C}} v_k \to f$ in $L^1(\gamma_n;\mathbb{C})$, i.e.
$\|f - \iota_{\mathbb{C}} v_k\|_{L^1} \to 0$.

\emph{Step 2: the two maps agree on the approximants.}
By clause (1) for $T$ and for $T'$,
$T(\iota_{\mathbb{C}} v_k) = R_\gamma^{\mathbb{C}} v_k
 = T'(\iota_{\mathbb{C}} v_k)$ for every $k$.  Writing
$f = (f - \iota_{\mathbb{C}} v_k) + \iota_{\mathbb{C}} v_k$ and using
$\mathbb{C}$-linearity of $T$ and of $T'$,
\[
  Tf - T'f = T\bigl(f - \iota_{\mathbb{C}} v_k\bigr)
    - T'\bigl(f - \iota_{\mathbb{C}} v_k\bigr) \qquad\text{for every } k .
\]

\emph{Step 3: the level sets.}
Let $\varepsilon > 0$.  If $|Tf(x) - T'f(x)|_{\mathbb{C}^n} > 2\varepsilon$ then,
by Step 2 and the triangle inequality in $\mathbb{C}^n$, at least one of
$|T(f - \iota_{\mathbb{C}} v_k)(x)| > \varepsilon$ and
$|T'(f - \iota_{\mathbb{C}} v_k)(x)| > \varepsilon$ holds.  Hence, by
monotonicity and subadditivity of $\gamma_n$ (shared context) and then clause
(2) applied to each of $T$ and $T'$ with the input $f - \iota_{\mathbb{C}} v_k$
and the level $\varepsilon$,
\[
  \gamma_n\bigl(\{|Tf - T'f| > 2\varepsilon\}\bigr)
    \leq \frac{8}{\pi\varepsilon}\|f - \iota_{\mathbb{C}} v_k\|_{L^1}
       + \frac{8}{\pi\varepsilon}\|f - \iota_{\mathbb{C}} v_k\|_{L^1}
    = \frac{16}{\pi\varepsilon}\,\|f - \iota_{\mathbb{C}} v_k\|_{L^1} .
\]

\emph{Step 4: conclusion.}
The left-hand side of the last display does not depend on $k$, while by Step 1
the right-hand side tends to $0$ as $k \to \infty$; a nonnegative constant
bounded by a sequence tending to $0$ is $0$, so
$\gamma_n(\{|Tf - T'f| > 2\varepsilon\}) = 0$ for every $\varepsilon > 0$.
Taking $\varepsilon = 1/j$ for $j \geq 1$ and using
\[
  \{Tf \neq T'f\} = \bigcup_{j \geq 1}\{|Tf - T'f| > 2/j\} ,
\]
a countable union of $\gamma_n$-null sets is null by subadditivity, so
$Tf = T'f$ $\gamma_n$-almost everywhere, i.e.\ $Tf = T'f$ in
$\mathcal{M}(\gamma_n;\mathbb{C}^n)$.  As $f \in L^1(\gamma_n;\mathbb{C})$ was
arbitrary, $T = T'$.
\end{proof}
